\songtitle{Este mar es azul sin disculpa}

\tag*{2026}\tag{milonga-rock}\tag{thomas-entourage}\tag{susana-bianchi}\tag{spanish-lyrics}
\begin{notes}
Lyrics by SUNO, based on Susana Bianchi's monologue in the \emph{Thomas’ Entourage} narrative universe by Carlos Thompson.
\end{notes}
\begin{style}
Milonga rioplatense contemporánea con textura de rock alternativo: compás de milonga marcado por guitarra criolla y bajo acústico, batería seca con escobillas que crece hacia golpes de rock contenidos, guitarras eléctricas con overdrive áspero y frases cortantes, bandoneón discreto en respuestas instrumentales, Voz femenina argentina, porteña y conversada, con dicción clara, fraseo narrativo y tensión íntima; versos casi recitados pero afinados, estribillos más abiertos y firmes, Arreglo en ascenso: comienzo desnudo, suma percusión y bajo en el segundo verso, guitarras amplias en los estribillos, puente instrumental con bandoneón y guitarra eléctrica, final contenido, Mezcla cálida y cercana, sala pequeña, graves definidos, saturación analógica moderada, sin brillo excesivo
\end{style}
\begin{lyrics}
\songpart{Intro}
El mar es el mismo desde cualquier lado,\\
eso lo sé, lo vi.\\
Punta Arenas, Dubái,\\
un cuarto en Fort Lauderdale,\\
Gregor dormido detrás de mí.

\songpart{Verse 1}
Yo miraba el Atlántico,\\
preguntándome qué hacía\\
tan lejos de casa,\\
tan lejos de la vida\\
que decía conocer.\\
El mar no tiene la culpa\\
de dónde le toca estar,\\
pero este mar, este mar,\\
me vuelve a preguntar.

\songpart{Chorus}
Este mar es azul sin disculpas,\\
horizonte limpio, filo de sal.\\
Si lo tuviera en el panel,\\
si pudiera elegir la vista,\\
quizás me quedaría acá.

\songpart{Verse 2}
El apartamento está bien,\\
y lo digo con honestidad:\\
anda el aire acondicionado,\\
la cama sabe descansar.\\
No hay mañanas con frazadas,\\
ni acuerdos con el clima;\\
acá el entorno hace lo que pido,\\
y eso mismo me lastima.

\songpart{Chorus}
Este mar es azul sin disculpas,\\
horizonte limpio, filo de sal.\\
Si lo tuviera en el panel,\\
si pudiera elegir la vista,\\
quizás me quedaría acá.

\songpart{Verse 3}
Dubái también funcionaba,\\
casi daba miedo mirar:\\
eficiencia sin fricción,\\
una ciudad sin accidente,\\
sin historia para tropezar.\\
Desierto atrás, mar adelante,\\
y en el medio, precisión;\\
una cosa enorme y perfecta\\
sin deberle nada a nadie, no.

\songpart{Bridge}
Me doy vuelta hacia la ventana,\\
la del oeste, la de las grúas.\\
Tres contra el cielo de la tarde,\\
dos quietas, una que continúa.\\
Y algo en ese movimiento lento\\
me calma más de lo que quisiera:\\
esto no es el resultado,\\
es apenas la manera.

\songpart{Verse 4}
Conozco planos y promesas,\\
corredores verdes por venir,\\
barrios unidos por la escala,\\
la ciudad aprendiendo a latir.\\
Natalia habla de integrar,\\
Thomas habla sin levantar la voz;\\
seriedad para lo que importa,\\
y yo escucho entre los dos.

\songpart{Final Chorus}
Este mar es azul sin disculpas,\\
pero no me puede decidir.\\
La grúa sigue contra la tarde,\\
la obra todavía respira,\\
y yo todavía estoy por ir.

\songpart{Outro}
El mar es el mismo desde cualquier lado,\\
pero no es el mismo para mí.\\
Este mar, este.\\
El que me encuentra lejos de casa,\\
y no me deja mentir.
\end{lyrics}

